\chapter{Conclusion}
\label{ref:conclusion}

\section{From automated layout generation to large-print adaptation}

This thesis argued that \melody's print production workflow could be automatized through intelligent, computational methods. This mattered given its dual impact: industrially, automation relieved the economic pressure faced by the press industry by reducing the cost of manual page composition; societally, the same spatial-layout techniques could be repurposed to serve readers under visual limitations. 

Automating this workflow unfolded along three axes:
\begin{itemize}
    \item Retrieving an appropriate template for each article, in both a single-article and a page-aware context, prior to composing the complete page.
    \item The central objective of this thesis: composing a complete page from a set of web articles without manual template assignment or placement, in a few minutes per page against the fifteen to twenty minutes currently required, while keeping objective values within the ranges observed in operator-produced pages.
    \item Quantifying the impact of large-print layouts for readers under constrained visual access, as a way to validate the re-layouting principle stated above.
\end{itemize}
Each of these axes drew on a distinct scientific discipline, with its own vocabulary, methodology, and evaluation criteria. This multidisciplinary breadth has been a defining characteristic of this thesis.

The following section summarizes the resulting three contributions, before turning to the perspectives they open.

\section{Main contributions and findings}
This thesis automatized \melody's web-to-print workflow by combining learned template retrieval and multi-objective page generation, and further evaluated the same underlying re-layouting principle for its applicability to newspaper large-print adaptation.

\subsection*{Contribution 1: Learned similarity-based template retrieval and page-aware scoring}
Template retrieval was addressed through the Template Retrieval System (\trs), which recommends templates to editors based on a learned structural similarity measure between templates, computed with a graph matching network and denoted $\gmn{\cdot}{\cdot}$, rather than by usage frequency alone. Relative to baselines, \trs was shown to nearly double structural diversity and to improve long-tail template exposure by approximately $1.5\times$, at a bounded loss of conformity to \melody's historical assignments, and at millisecond-range latency.

This similarity-based scoring was further extended to a page-aware setting through the Geometric-Weighted Editorial Score (\neuralmethod), which combines a Transformer encoder, reaching $82.62\%$ Top-1 accuracy in predicting \melody's historical assignments, with a matrix of pairwise structural similarity scores $\gmnMatrix$, computed with the same $\gmn{\cdot}{\cdot}$ metric introduced above. This expanded the average pool of viable template alternatives per article by $1.7\times$, so as to assist the search performed by the multi-objective page generation method described next, which incorporates this scoring as one of its optimization objectives.

\begin{graybox}
\trs was deployed on \melody's pre-production infrastructure to be used by \publihebdos press group. \neuralmethod's deployment is reported as part of the next contribution, since it is folded into \methodNameGeneral.
\end{graybox}

\subsection*{Contribution 2: Multi-objective page generation}

Automating the page generation objective was addressed through \methodNameGeneral, a multi-objective method that jointly resolves template assignment and two-dimensional packing across four editorial objectives. Its two variants, uniform \methodNameUniform and preference-guided \methodName, were shown to discover Pareto-dominating solutions over constructive and ablated baselines across the large majority of evaluated instances.
\methodName further compressed the non-dominated frontier by an order of magnitude and concentrated up to $51.4\%$ more solutions within the acceptable range of historical operator-produced pages.

\begin{graybox}
\methodNameGeneral, which incorporates \neuralmethod's scoring as one of its optimization objectives, was deployed on \melody's pre-production infrastructure to be used by \laprovence press group.
\end{graybox}

\subsection*{Contribution 3: Large-print validation under constrained visual access}

Under the same automated re-layouting principle, large-print adaptation addressed the societal motivation of this thesis: validating that large-print layouts reduce reading time and perceived workload, and better preserve the behavioral integrity of reading, relative to gesture-based magnification, for readers \NV under device-related visual constraints and for readers \LV.

A controlled user study showed that \condtwonot restores the behavioral integrity of the reading process, with reading paths remaining significantly closer to the natural, left-to-right and top-to-bottom strategy than under \condonenot, while also yielding faster task completion and lower perceived workload. In conclusion, \condtwonot can therefore be regarded as an instance of inclusive design~\cite{langdon-inclusive-design-2014}. \ \\

Together, these three contributions substantiate the argument that \melody's print production workflow can be automatized through intelligent, computational methods, yielding \textbf{both industrial gains in production efficiency and a societal extension toward readability}.


\section{Perspectives}
\label{sec:perspectives}

Each of the three contributions summarized above considers one of the scientific challenges stated in \secref{sec:scientific-challenges-contributions}, \textbf{and each one also delimits a boundary that this thesis did not cross}. The perspectives developed below follow those boundaries outwards: from extending the automated workflow to the complete newspaper, to adapting that layout to individual readers, to admitting human judgment into the optimization process itself, and finally to questioning the fixed two-dimensional page that every formulation in this thesis has taken for granted.

\subsection{Generalizing the automated workflow}
\label{sec:perspectives-workflow}

\textbf{A broader research question raised by this thesis concerns how newspaper layout generation can be extended from the automated construction of individual pages to the end-to-end generation of complete newspaper editions.} \methodNameGeneral\ addresses several components of this process, but still treats page generation as an isolated problem: the set of articles to be composed is given, and the canvas to be filled is a single page. Producing a complete edition instead requires jointly addressing other decisions, from selecting and assigning content to pages, to determining their layouts and placing every element they carry, including classes of element that this thesis left out entirely. 

The clearest such omission concerns advertisements, which the current implementation of \methodNameGeneral\ does not place. This omission stems from a data limitation rather than from the formulation itself: the historical layout records available at \melody\ lack consistent metadata describing advertisement coordinates and dimensions, leaving no reference against which their placement could be calibrated or evaluated. The current formulation already can accommodate them as an additional class of placeable object, as detailed in \secref{sec:pagegen-discussion}. Since their aspect ratios are predetermined, advertisements can be modeled as static, fixed-size bounding boxes that bypass the template selection and feed directly into the page placement. On the other hand, a spatial penalty, formulated analogously to the heading critical zone penalty $\HCZ$, steers them towards the quadrants reserved for them. This extension is therefore restricted by the availability of the corresponding metadata rather than by any further methodological development.

A second direction concerns the single-page restriction itself. The positional variable of the formulation already carries a page index $\pageindex{i}$, which the feasibility constraint FC~1 fixes to a single canvas for every article (\secref{sec:constraints_space}). 
Relaxing that constraint would nonetheless not be sufficient to produce a complete edition, since generating one additionally requires assigning each article to one of several pages according to its editorial section, such as sports or politics, and to the overall structure of the edition. 
This constitutes a third assignment problem, layered on top of the template assignment and two-dimensional packing problems that \methodNameGeneral\ already resolves jointly: the section assigned to an article determines which pages may receive it, which in turn constrains both the templates that remain admissible for it and the space left to pack them. Scaling from a single page to a complete edition therefore alters the nature of the optimization problem.

A last direction concerns generalization to other contexts. All empirical work in this thesis is grounded in a single CMS, \melody, and validated with two of its press groups, \publihebdos\ and \laprovence, both operating within the French regional press. The problem does not depend on that context: neither the representations nor the optimization mechanisms developed here are formulated in terms of a particular publisher. However, specific parameters  were imposed by \melody's experts rather than derived from user studies, and whether the four editorial objectives and the coefficients weighting them would hold in another market, language, or editorial culture remains an open question. 
Answering this would require evaluating the proposed approaches across different CMSs, publishers, markets, and languages, or in structurally related layout domains such as magazines or user interfaces. \
\\
Taken together, these three directions shift the problem from automating the construction of a page to automating the production of a complete editorial layout, while asking how such an approach can remain adaptable to different content, editorial structures, and production environments.

\subsection{From large-print adaptation to personalized layouts}
\label{sec:perspectives-personalization}

The validation reported in Chapter~\ref{chapter:magnification} established the benefit of large-print layouts. However, the stimuli evaluated there were synthesized by an earlier, entry-point-preserving re-layouting method~\citemine{gallardo} that predates \methodNameGeneral\ and operates on text-only pages, enlarging headlines and body text while preserving a two-dimensional layout format. \methodNameGeneral\ itself, which composes a complete page from a full set of web articles across a template catalog, has not been adapted to generate complete large-print editions, incorporating images and other non-textual elements alongside enlarged text. A gap therefore separates the principle that was validated from the exact method that was deployed.

Closing that gap is primarily a matter of templates rather than of the optimization method. The catalog \templatecatalog\ available at \melody\ was designed for standard print sizes: its zoning topologies, and the text capacity matrices $\capacity{j}$ derived from them, presuppose body text at its standard size, and leave correspondingly little room for a magnified font size. Producing complete large-print editions therefore requires first deriving alternative, large-print variants of the templates themselves, with fewer and larger zones and recomputed capacities. How such variants should be obtained, whether by transforming existing templates or by designing them from scratch, is left open. 
Once such a catalog is available, \methodNameGeneral\ applies without structural modification, since neither its search nor its packing machinery assumes a particular print size. What must change is the criterion by which it selects among templates: the editorial-style objective $\objtwo$ currently rewards conformity with \melody's historical assignments, as learned by the Transformer encoder of \neuralmethod, whereas a large-print edition would require it to reward readability, for which no comparable historical corpus exists. The remaining effort thus lies in engineering the template catalog, and not in reformulating but adapting the page generation method.

A further question concerns the adaptation itself. The large-print condition evaluated in this thesis applies one common adaptation to every reader. This raises a broader research question: \textbf{how can newspaper layouts adapt to individual visual and reading needs, drawing on a reader's measured parameters, expressed preferences, and observed reading behavior?}

Mechanisms of this kind already exist for plain text, where \emph{SituFont}~\cite{yue2024situfont} adjusts typographic attributes from sensor input and human-in-the-loop feedback, and where gaze-driven systems augment the reading experience in real time~\cite{aguilar-castet2017, wang_gazeprompt_2024}; what remains open is their extension to the spatial organization of a layout. 
Pursuing this direction would also require extending the empirical validation reported here, since the reading-aloud measure captures only textual entry points (\secref{sec:magnification-discussion}).

Structuring such studies around a participatory design methodology~\cite{schuler-participatory-design-1993, wacnik-participatory-review-2024}, involving readers \NV\ and \LV\ throughout the design process rather than solely at the evaluation stage, would establish inclusive design~\cite{langdon-inclusive-design-2014} as an explicit priority of this line of work rather than an auxiliary consideration.

\subsection{Towards human-guided layout optimization}
\label{sec:perspectives-human-guided}

The participatory framing mentioned above applies equally to the evolutionary framework itself. The editorial objectives optimized by \methodNameGeneral, together with the domain-specific coefficients weighting them, were calibrated through focus-group meetings with \melody\ experts, which constitutes a one-off consultation rather than a sustained involvement of those experts in the design of the method. Once objectives and weights were fixed, moreover, the architecture offered no mechanism for an expert to act on an individual solution, for instance by adapting a proposed page to their own judgment after it had been generated. Both the definition of the editorial objectives and the evaluation of the resulting layouts therefore rest on human judgment, yet that judgment remains external to the optimization process. \textbf{A key research question is accordingly how expert and reader feedback could become part of the optimization loop itself.}

Such a mechanism is already latent in the formulation. Preferences can be articulated at three points relative to a multi-objective search: a priori, by supplying target values before it begins; interactively, by refining the search after inspecting an initial approximation of the front; or a posteriori, by selecting a solution from a fully generated front~\cite{rnsga-iii}. \methodName\ exploits the first of these, deriving its aspiration points from historical layouts, which matched \melody's immediate aim of generating pages automatically as a first step. Admitting the expert into the evolutionary search itself is the natural next one~\cite{takagi-iec-2001, afsar-interactive-moo-2021}: through the interactive mode, an editor would inspect a first frontier and reposition the aspiration points so that the following iteration concentrates solutions in the region actually sought. This is a genuine development rather than the activation of an existing option, since it requires an interaction model, a revision protocol, and a computational budget compatible with the response times expected in such a CMS. It would nonetheless replace the one-off calibration of the editorial objectives by a continuous one, and open the way to validating them against a wider panel of operators or end readers, rather than against the judgment of a limited group of experts.

\subsection{From fixed pages to spatial and wearable media}
\label{sec:perspectives-spatial}

Every formulation in this thesis presupposes a page: a fixed, rectangular, two-dimensional surface onto which content is to be arranged. \textbf{What was never questioned is the two-dimensionality itself.} 
The newspaper page has historically been shaped by the physical constraints of its medium, and the transition from print to desktop displays, tablets and smartphones has progressively relaxed some of those constraints while introducing others. Emerging spatial and wearable media may represent a further shift: rather than rendering a page onto a fixed rectangular surface, they could allow editorial content to be organized dynamically within the reader's visual and physical space. The question addressed throughout this thesis, namely how content should be arranged on a given page, would then give way to a different one: given a content, a reader, and a medium, how should the reading space itself be constructed?

Answering such a question would require rethinking editorial hierarchy, readability, navigation and accessibility for a medium that none of these notions was formulated for. \textbf{Where this thesis drew on a multidisciplinary schema, this last perspective would draw on at least four, none of which it engaged with.} The medium itself belongs to \emph{virtual and augmented reality} and to \emph{wearable displays}, whose editorial use is beginning to be examined under the heading of \emph{immersive journalism}~\cite{bujic-imx2023}.
Readability belongs to \emph{the psychophysics of reading}, already medium-independent in the relevant respect, since print size is defined as a visual angle rather than as a physical dimension~\cite{legge2006} and therefore remains meaningful in a space that contains no page. 
Navigation belongs to \emph{human-computer interaction}, whose established principles were formulated for flat screens: the progression from overview to zoom and filter, and then to details on demand~\cite{shneiderman1996}, presupposes a bounded display that a three-dimensional reading space does not provide.
And accessibility belongs to \emph{inclusive and participatory design}~\cite{langdon-inclusive-design-2014, schuler-participatory-design-1993, vatavu-imx2021}, which would require readers \LV\ to take part in designing such media from the start, and not once they are already built.

This direction remains entirely speculative, and has not been formalized beyond the framing given here. It is nonetheless an interesting one, and of potentially high impact for the future of the field, where the representations and optimization tools developed in this thesis may prove useful.

\bigskip

\noindent Taken together, these four directions trace a single trajectory, along which each one relaxes an assumption that the preceding work relied on: the page gives way to the edition, a common adaptation to a per-reader one, a fixed calibration to a continuous interaction, and a bounded rectangular surface to a reading space that must itself be constructed. What remains constant across them is the premise defended throughout this thesis, namely that the spatial organization of editorial content is a computational problem, and that automating it serves both the production of newspapers and the readers who navigate them.